\section{\genregating: the (domain $\times$ genre) switch}
\label{sec:genre}

\begin{figure}[t]
\centering
\includegraphics[width=\linewidth]{fig6_genre.pdf}
\caption{Genre gating across two replicated measurements. (i) Repair-component
gains are visible in the repair genre ($+10$--$16$pp over placebo, three seeds
per arm) but not on plain items ($+2$--$7$pp). (ii) Drills harm at a 25\% dose
appears only in the repair genre (single preregistered dose arm; the pure-diet
endpoint is two-seed). A retired third case is autopsied in
Appendix~\ref{app:retired}.}
\label{fig:genre}
\end{figure}

Why should the evaluation genre matter at all? Because the model's failures
are themselves genre-coupled: the Janet item of Table~\ref{tab:probes} is
answered correctly as prose and incorrectly (\texttt{\{"answer": 26\}}) as
JSON, by the same model on the same arithmetic. What holds for single probes
holds for whole evaluations. The same checkpoints tell different stories
depending on the genre in which they are asked to perform, and the pattern recurs across independent, multi-seed
measurements (Figure~\ref{fig:genre}). Repair-component gains that are barely
visible on plain items ($+2$--$7$pp over placebo) separate cleanly in the
repair genre ($+10$--$16$pp). Drills toxicity is invisible on plain items up
to a 25\% dose yet already collapses repair-genre keep accuracy by 17pp at
that dose. A third GSM case --- a sign flip in the scaffold component's effect
between genres --- was observed on a single seed at a marginal-health operating
point and did not survive seed replication (2 of 3 seeds reverse at healthy
operating points); it is retired with an autopsy (Appendix~\ref{app:retired}). Of the cases that
remain, gain visibility rests on three seeds per arm and the 2Wiki reversal on
three seeds; the dose-level drills case is a single preregistered arm whose
pure-diet endpoint is two-seed.

\begin{figure}[t]
\centering
\includegraphics[width=\linewidth]{fig8_validation.pdf}
\caption{The pipeline off GSM. (a--b) 2WikiMultihopQA miniature: the
anti-credulity component transfers on the plain face (adopt 19$\to$0) while
repair-genre accuracy does not separate from placebo and every 600-item arm
sits below the \prerepair{}. (c) 200 real-world messy prompts from six public
datasets map onto the probe taxonomy; the blind-spot column (retrieval noise)
is reported honestly. Operating points e4/e2; one representative seed shown, three-seed aggregates in Table~\ref{tab:wiki3seed}.}
\label{fig:validation}
\end{figure}

The validation domain shows the gate is a genuine switch, not a one-way
amplifier. On 2WikiMultihopQA (Figure~\ref{fig:validation}a--b) the direction
reverses: the anti-credulity component transfers on the plain face --- adoption
failures drop from 19 to 0 and derailments from 28 to 4 against a placebo at
7/56 --- while repair-genre accuracy does not separate ($+3.8$pp, below the
preregistered $+8$pp line, scored as a miss). Which genre reveals an effect is
domain-dependent; what is constant is that no single genre suffices. The 2Wiki
miniature also surfaced an unregistered cost specific to the knowledge domain:
every 600-item arm, including the placebo, damages known-item QA (5
$\rightarrow$ 34--56 new failures), and a transcript audit classifies 100\% of
those as answered-wrong with confabulated evidence chains --- a content-side
tax, distinct from the genre-bleed tax of the computation domain --- and every
arm lands below the \prerepair{} on repair-genre accuracy. The prescription
consequence is rule 2 of \S\ref{sec:prescription}: validate in both genres,
against the \prerepair{}, without assuming a priori where the effect will show.
